% SPDX-License-Identifier: PMPL-1.0-or-later
% Copyright (c) 2026 Jonathan D.A. Jewell (hyperpolymath) <j.d.a.jewell@open.ac.uk>
%
% Gossamer: A Linearly-Typed Webview Shell with Provable Resource Safety
% Prepared for submission to arXiv (cs.PL / cs.SE)

\documentclass[11pt,a4paper]{article}

\usepackage[utf8]{inputenc}
\usepackage[T1]{fontenc}
\usepackage{amsmath,amssymb,amsthm}
\usepackage{mathpartir}
\usepackage{stmaryrd}
\usepackage{booktabs}
\usepackage{multirow}
\usepackage{graphicx}
\usepackage{listings}
\usepackage{xcolor}
\usepackage{hyperref}
\usepackage[margin=2.5cm]{geometry}
\usepackage{enumitem}
\usepackage{caption}
\usepackage{subcaption}

% Code listing style
\definecolor{codebg}{rgb}{0.96,0.96,0.96}
\definecolor{codegreen}{rgb}{0.0,0.5,0.0}
\definecolor{codepurple}{rgb}{0.5,0.0,0.5}
\definecolor{codeorange}{rgb}{0.8,0.4,0.0}

\lstdefinestyle{ephapax}{
  backgroundcolor=\color{codebg},
  basicstyle=\ttfamily\small,
  keywordstyle=\color{codepurple}\bfseries,
  commentstyle=\color{codegreen}\itshape,
  stringstyle=\color{codeorange},
  morekeywords={fn,let,let!,region,match,import,type,module,in,if,else},
  breaklines=true,
  frame=single,
  framerule=0.5pt,
  numbers=left,
  numberstyle=\tiny\color{gray},
  xleftmargin=2em,
}

\lstdefinestyle{idris}{
  backgroundcolor=\color{codebg},
  basicstyle=\ttfamily\small,
  keywordstyle=\color{codepurple}\bfseries,
  commentstyle=\color{codegreen}\itshape,
  morekeywords={module,import,data,where,public,export,Type,IO,Either,Maybe,So,Bits64,Bits32},
  breaklines=true,
  frame=single,
  framerule=0.5pt,
  numbers=left,
  numberstyle=\tiny\color{gray},
  xleftmargin=2em,
}

\lstdefinestyle{zig}{
  backgroundcolor=\color{codebg},
  basicstyle=\ttfamily\small,
  keywordstyle=\color{codepurple}\bfseries,
  commentstyle=\color{codegreen}\itshape,
  morekeywords={export,fn,const,var,pub,switch,orelse,return,null,void,struct,enum,import,catch,defer},
  breaklines=true,
  frame=single,
  framerule=0.5pt,
  numbers=left,
  numberstyle=\tiny\color{gray},
  xleftmargin=2em,
}

\theoremstyle{definition}
\newtheorem{definition}{Definition}[section]
\newtheorem{theorem}{Theorem}[section]
\newtheorem{lemma}[theorem]{Lemma}
\newtheorem{corollary}[theorem]{Corollary}
\newtheorem{property}{Property}[section]

\title{Gossamer: A Linearly-Typed Webview Shell\\with Provable Resource Safety}

\author{
  Jonathan D.A. Jewell\\
  \texttt{j.d.a.jewell@open.ac.uk}\\
  hyperpolymath
}

\date{March 2026}

\begin{document}

\maketitle

\begin{abstract}
Desktop application frameworks that wrap web frontends in native webview
windows --- Electron, Tauri, Wails, and others --- universally manage
resources through garbage collection, reference counting, or manual
tracking, and communicate across the frontend--backend boundary via
untyped JSON serialisation. We present \textsc{Gossamer}, a webview
shell framework built in Ephapax, a dyadic language with both affine
and linear type modes. Gossamer introduces three novel contributions:
(1)~linearly-owned webview handles with compile-time lifecycle proofs,
(2)~a typed IPC protocol where frontend--backend message shape agreement
is verified at compile time rather than at runtime, and
(3)~linear capability tokens that enforce access control through the
type system rather than through runtime configuration checks.
We formalise the core type system, describe the implementation
architecture using Idris2 ABI definitions and Zig FFI bindings, and
present a comprehensive comparison with fourteen existing frameworks
across safety, performance, and developer experience dimensions.
Gossamer is, to our knowledge, the first webview shell where the type
system proves that resources cannot leak, handles cannot be used after
destruction, and permissions cannot be bypassed.
\end{abstract}

\section{Introduction}
\label{sec:intro}

The webview shell pattern --- a native window containing an embedded
web rendering engine, connected to a system-level backend via
inter-process communication --- has become the dominant architecture
for cross-platform desktop applications. From Slack and Discord
(Electron) to emerging projects built on Tauri, developers choose this
pattern because it allows reuse of web frontend skills and codebases
while providing access to native operating system capabilities.

Despite this dominance, all existing implementations share fundamental
weaknesses rooted in how they manage two concerns: \emph{resource
lifecycle} and \emph{type safety at the IPC boundary}.

\paragraph{Resource lifecycle.} A webview shell manages several kinds
of resources: the webview handle itself (which must be created on the
main thread and properly destroyed), IPC channels, event listener
registrations, plugin-held resources such as file handles or database
connections, and system tray icons. Existing frameworks manage these
through garbage collection (Electron, Wails, Flutter), reference
counting (Tauri's \texttt{Arc<Mutex<...>>}), or manual tracking
(CEF, Neutralinojs). All three approaches admit resource leaks,
use-after-free, and double-free as runtime failure modes.

\paragraph{IPC type safety.} The frontier between frontend (JavaScript)
and backend (Rust, Go, C++, Node.js) is universally a JSON serialisation
boundary. Electron passes JavaScript objects through
\texttt{contextBridge}; Tauri serialises Rust structs via Serde and
deserialises on the frontend; Wails uses Go reflection to bind struct
methods to JavaScript. In all cases, a mismatch between the shape
expected by the frontend and the shape produced by the backend manifests
as a runtime error --- a \texttt{JSON.parse} failure, an
\texttt{undefined} field access, or a silent data truncation.

\paragraph{Permission enforcement.} Tauri introduced a
capability-based permission system where each window or plugin must
declare which system APIs it needs access to. This is the most
sophisticated security model in the field, but it operates at runtime:
a JSON configuration file is parsed at startup, and each API call is
checked against it. A misconfigured file silently permits unintended
access; the compiler cannot detect the error.

\medskip

We propose \textsc{Gossamer}, a webview shell framework that addresses
all three weaknesses by building on Ephapax, a dyadic programming
language with both affine and linear type modes, region-based memory
management, and a formally verified ABI layer defined in Idris2.

\section{Background}
\label{sec:background}

\subsection{Ephapax: A Dyadic Linear Language}

Ephapax~\cite{ephapax2025} is a programming language with a dyadic type
system offering two modes of substructural typing:

\begin{itemize}[nosep]
  \item \textbf{Affine mode}: values may be used \emph{at most once}.
    Unused values are implicitly dropped (cleanup at region exit).
  \item \textbf{Linear mode}: values must be used \emph{exactly once}.
    The compiler rejects any program where a linear value is unused or
    used more than once.
\end{itemize}

Both modes share the same FFI surface and the same compiled
representation; the distinction is enforced entirely at compile time.
The key syntactic marker is the \texttt{let!} binding for linear values
versus \texttt{let} for affine values.

Ephapax targets WebAssembly via a two-phase compiler: an Idris2 stage
produces S-expression intermediate representation, which a Rust backend
compiles to WASM. The \texttt{ephapaxiser} tool provides linearity
enforcement for existing codebases.

\subsection{Region-Based Memory Management}

Ephapax uses Tofte--Talpin style region-based memory
management~\cite{tofte1997}. Allocations are scoped to a named region;
when the region exits, all allocations within it are freed in bulk,
with zero per-object overhead:

\begin{lstlisting}[style=ephapax]
region r:
    let s1 = String.new@r("hello")
    let len = String.len(&s1)
    let result = String.concat(s1, s2)
    result
// region r exits: all allocations freed
\end{lstlisting}

The classical weakness of region-based systems is \emph{region escape}:
a value allocated in region $R$ being returned to a scope where $R$
has been deallocated. Existing systems handle this through conservative
region inference (MLKit~\cite{tofte1997}), existential types with
runtime checks (Cyclone~\cite{jim2002cyclone}), or reference counting
as a fallback (Rust's \texttt{Arc}/\texttt{Rc}).

Ephapax's linear types close this escape hatch: a linear value bound
within a region must be consumed within that region's scope. The
compiler statically rejects any attempt to return a linear value past
its region boundary.

\begin{theorem}[No Region Escape]
\label{thm:no-escape}
In a well-typed Ephapax program, a value allocated in region $r$
cannot appear in any expression that outlives $r$. This holds for
\emph{both} affine and linear bindings.
\end{theorem}

\begin{proof}
The type checker's \texttt{check\_region} method enforces
\textsc{NoRegionInType}: before returning from a region block, the
return type is checked via \texttt{Ty::references\_region(r)}, which
recursively inspects all type constructors (String, Region, Ref, Fun,
Prod, Sum, List, Tuple, Borrow). If the return type references $r$,
the program is rejected. This check is \emph{qualifier-independent}:
it inspects the type, not the binding's qualifier. Therefore it applies
identically to affine (\texttt{let}) and linear (\texttt{let!})
bindings. \qed
\end{proof}

\begin{theorem}[All Linears Consumed at Region Exit]
\label{thm:all-linears-consumed}
At region exit, all linear variables bound within the region have been
consumed. Affine variables may be implicitly dropped.
\end{theorem}

\begin{proof}
The type checker tracks which region each variable was bound in
(\texttt{CtxEntry.region}). At region exit, it iterates all variables
with \texttt{region = Some(r)} and verifies that every linear variable
(\texttt{ty.is\_linear() == true}) has been marked as used. Affine
variables are not checked---their memory is freed by the region's
arena deallocator. \qed
\end{proof}

\begin{lemma}[Qualifier-Region Orthogonality]
\label{lem:orthogonality}
The region escape check (\Cref{thm:no-escape}) and the qualifier
enforcement (linear must consume, affine may drop) are independent.
Changing a binding's qualifier does not affect whether it can escape.
Changing a binding's region does not affect whether it must be consumed.
\end{lemma}

\begin{proof}
By inspection of the type checker implementation:
\texttt{Ty::references\_region} does not inspect the binding's
qualifier; the \texttt{Split} rules do not inspect the binding's
region. The two concerns compose without interaction. \qed
\end{proof}

\subsection{The Idris2 ABI / Zig FFI Standard}

The hyperpolymath ABI/FFI standard~\cite{hyperpolymath-abi} defines
a three-layer architecture for foreign function interfaces:

\begin{enumerate}[nosep]
  \item \textbf{Idris2 ABI layer}: interface definitions with dependent
    type proofs (memory layout, handle validity, protocol shapes).
  \item \textbf{Zig FFI layer}: C-compatible implementations with
    cross-compilation support and zero runtime dependencies.
  \item \textbf{Generated C headers}: bridge between ABI and FFI.
\end{enumerate}

Idris2's dependent types enable formal proofs of ABI properties that
are impossible to express in C, Zig, or Rust alone. For example, the
\texttt{Handle} type carries a non-null proof:

\begin{lstlisting}[style=idris]
data Handle : Type where
  MkHandle : (ptr : Bits64)
           -> {auto 0 nonNull : So (ptr /= 0)}
           -> Handle
\end{lstlisting}

\section{Gossamer Architecture}
\label{sec:architecture}

\subsection{Overview}

Gossamer is structured as four layers:

\begin{enumerate}[nosep]
  \item \textbf{Web frontend}: standard HTML/CSS/JS (or ReScript)
    running inside the OS-provided webview.
  \item \textbf{Ephapax application layer}: command handlers with linear
    ownership, capability tokens, and the plugin host.
  \item \textbf{Idris2 ABI layer}: formal specifications of the webview
    handle, IPC protocol, and capability types.
  \item \textbf{Zig FFI layer}: platform-specific bindings to
    WebKitGTK (Linux), WKWebView (macOS), and WebView2 (Windows).
\end{enumerate}

\subsection{Linearly-Owned Webview Handle}
\label{sec:webview}

The webview handle is modelled as a linear resource with a dependent
type proof of non-nullity and a main-thread witness:

\begin{lstlisting}[style=idris]
data WebviewHandle : Type where
  MkWebview : (ptr : Bits64)
            -> {auto 0 nonNull : So (ptr /= 0)}
            -> {auto 0 onMainThread : MainThreadProof}
            -> WebviewHandle
\end{lstlisting}

Operations on the handle fall into two categories:

\begin{itemize}[nosep]
  \item \textbf{Borrowing operations} (\texttt{loadHTML},
    \texttt{navigate}, \texttt{eval}): accept the handle linearly and
    return it alongside the operation result. The caller retains
    ownership.
  \item \textbf{Consuming operations} (\texttt{run}, \texttt{destroy}):
    accept the handle linearly and do not return it. The handle is
    irrecoverably consumed.
\end{itemize}

This encoding provides the following compile-time guarantees:

\begin{property}[No Use-After-Free]
After a consuming operation (\texttt{run} or \texttt{destroy}), the
handle variable is no longer in scope. Any subsequent reference is a
type error.
\end{property}

\begin{property}[No Double-Free]
Since the handle is linear, it can be passed to exactly one consuming
operation. Attempting to pass it to two is a linearity violation.
\end{property}

\begin{property}[No Leak]
A linear handle that is neither passed to a borrowing operation nor a
consuming operation before its region exits triggers a compile-time
``unused linear variable'' error.
\end{property}

\subsection{Typed IPC Protocol}
\label{sec:ipc}

The IPC channel is parameterised by request and response types:

\begin{lstlisting}[style=idris]
data Channel : (req : Type) -> (resp : Type) -> Type where
  MkChannel : (ptr : Bits64)
            -> {auto 0 nonNull : So (ptr /= 0)}
            -> Channel req resp
\end{lstlisting}

A protocol is a type-level list of named commands, each with typed
request and response:

\begin{lstlisting}[style=idris]
data Protocol : List (String, Type, Type) -> Type where
  Nil  : Protocol []
  (::) : (name : String, req : a, resp : b)
       -> Protocol rest
       -> Protocol ((name, a, b) :: rest)
\end{lstlisting}

Command handlers are registered with full type information:

\begin{lstlisting}[style=idris]
handle : (1 _ : Channel req resp)
       -> (name : String)
       -> (handler : req -> IO resp)
       -> IO (Channel req resp)
\end{lstlisting}

On the frontend side, a generated JavaScript client (produced from the
protocol type at compile time) provides a typed interface:

\begin{lstlisting}[language=JavaScript]
// Generated from Protocol [("greet", String, String)]
const gossamer = {
  greet: async (name: string): Promise<string> => {
    return __gossamer_invoke("greet", name);
  }
};
\end{lstlisting}

\begin{property}[IPC Agreement]
If the Ephapax backend type-checks and the JavaScript client is
generated from the same protocol definition, then every frontend
invocation matches a backend handler in name, request type, and
response type.
\end{property}

\subsection{Linear Capability Tokens}
\label{sec:capabilities}

Capabilities are modelled as linear values parameterised by resource
kind:

\begin{lstlisting}[style=idris]
data Cap : (resource : ResourceKind) -> Type where
  MkCap : (token : Bits64) -> Cap resource

data ResourceKind
  = FileSystem FilesystemScope
  | Network NetworkScope
  | Shell ShellScope
  | Clipboard
  | Notification
  | Tray
\end{lstlisting}

A plugin or window receives capability tokens at initialisation. The
linearity constraint ensures:

\begin{itemize}[nosep]
  \item \textbf{No forgery}: the \texttt{MkCap} constructor is not
    exported; capabilities can only be created by the framework.
  \item \textbf{No duplication}: linear values cannot be copied.
    A capability token is held by exactly one owner.
  \item \textbf{Revocability}: the \texttt{revoke} operation consumes
    the token, preventing further use.
  \item \textbf{Scope enforcement}: \texttt{FilesystemScope} restricts
    access to specific paths, checked at the type level.
\end{itemize}

\begin{property}[Capability Soundness]
A well-typed Gossamer program cannot perform a capability-gated
operation without holding a corresponding \texttt{Cap} value in its
linear context.
\end{property}

This contrasts with Tauri's approach, where capabilities are runtime
configuration. A Tauri application with a permissive
\texttt{tauri.conf.json} silently grants broad access. In Gossamer,
the capability must exist as a value in the program --- if it is not
threaded through to the operation site, the program does not compile.

\section{Implementation}
\label{sec:implementation}

\subsection{Zig FFI Layer}

The Zig FFI layer provides a platform-abstracted C ABI surface.
Platform selection uses compile-time dispatch:

\begin{lstlisting}[style=zig]
const impl = switch (@import("builtin").os.tag) {
    .linux   => @import("webview_gtk.zig"),
    .macos   => @import("webview_cocoa.zig"),
    .windows => @import("webview_win32.zig"),
    else     => @compileError("Unsupported platform"),
};

export fn gossamer_create(
    title: [*:0]const u8,
    width: u32, height: u32,
    resizable: u8, decorations: u8,
) ?*GossamerHandle {
    return impl.create(title, width, height, resizable, decorations);
}
\end{lstlisting}

The Linux implementation binds to WebKitGTK via \texttt{@cImport},
calling \texttt{gtk\_window\_new}, \texttt{webkit\_web\_view\_new},
and \texttt{gtk\_main} for the event loop. macOS (WKWebView) and
Windows (WebView2) stubs are provided for Phase~2.

\subsection{FFI Intrinsic and Native Execution}

Ephapax provides an \texttt{\_\_ffi} intrinsic that calls C ABI
functions from loaded shared libraries via \texttt{dlopen}/\texttt{dlsym}:

\begin{lstlisting}[style=ephapax]
fn main(): I64 =
  let! handle = __ffi("gossamer_create", "Hello", 800, 600, 1, 1, 0) in
  let! _ = __ffi("gossamer_load_html", handle, "<h1>Hello!</h1>") in
  __ffi("gossamer_run", handle)
\end{lstlisting}

The type checker allows string literals in FFI context without region
allocation---they become temporary C strings at the boundary. Linear
variables passed to FFI are consumed (counts as consumption for the
linearity check). The interpreter resolves symbols at runtime:

\begin{verbatim}
$ ephapax run hello.eph -L libgossamer.so
\end{verbatim}

This has been verified end-to-end: \texttt{gossamer\_create} returns a
real GTK window handle, \texttt{gossamer\_cap\_grant} returns
cryptographic tokens, and chained calls (\texttt{cap\_grant} $\to$
\texttt{cap\_check}) produce correct results through native execution.

\subsection{Ephapax Application Layer}

A minimal Gossamer application using the actual compiler syntax:

\begin{lstlisting}[style=ephapax]
fn main(): I64 =
  let! handle = __ffi("gossamer_create", "Hello Gossamer", 800, 600, 1, 1, 0) in
  let! _ = __ffi("gossamer_load_html", handle, "<html>...</html>") in
  __ffi("gossamer_run", handle)
\end{lstlisting}

The linear bindings (\texttt{let!}) enforce that:
\begin{itemize}[nosep]
  \item The handle returned by \texttt{gossamer\_create} \emph{must} be
    consumed (passed to \texttt{gossamer\_run} or \texttt{gossamer\_destroy}).
  \item Using \texttt{handle} after \texttt{gossamer\_run} is a type error
    (the variable has been consumed by the FFI call).
  \item Omitting the \texttt{gossamer\_run} call is a type error (linear
    variable not consumed).
\end{itemize}

This has been verified with the actual Ephapax compiler: the program
type-checks in linear mode and executes via native FFI, creating a real
GTK/WebKitGTK window.

\subsection{Dyadic FFI Design}

Following the Ephapax dyadic FFI pattern established in the
\texttt{proven} library bindings, affine and linear modes share a
single Zig FFI layer. The distinction is purely at the Ephapax
compiler level:

\begin{lstlisting}[style=ephapax]
// Affine mode — implicit cleanup OK
fn affine_example() -> i32 {
    region proto:
        let webview = Shell.create@proto({...})
        let webview = Shell.loadHTML(webview, "...")
        Shell.run(webview)
        // channel cleanup is implicit at region exit
        0
}

// Linear mode — explicit consumption required
fn linear_example() -> i32 {
    region prod:
        let! webview = Shell.create@prod({...})
        let! channel = Bridge.open@prod(webview)
        let! webview = Shell.loadHTML(webview, "...")
        Shell.run(webview)        // consumes webview
        Bridge.close(channel)     // REQUIRED
        0
}
\end{lstlisting}

\section{Comparison with Existing Frameworks}
\label{sec:comparison}

Table~\ref{tab:safety} compares safety properties across fourteen
existing frameworks and Gossamer.

\begin{table}[ht]
\centering
\caption{Safety and correctness comparison}
\label{tab:safety}
\small
\begin{tabular}{@{}lccccc@{}}
\toprule
\textbf{Framework} & \textbf{Resource} & \textbf{IPC Type} & \textbf{Permission} & \textbf{Memory} & \textbf{GC} \\
 & \textbf{Lifecycle} & \textbf{Safety} & \textbf{Model} & \textbf{Safety} & \textbf{Required?} \\
\midrule
Electron      & GC           & None       & Opt-in sandbox  & GC         & Yes \\
NW.js         & GC           & None       & None            & GC         & Yes \\
CEF           & Refcount     & None       & Process sandbox & Manual     & Partial \\
Tauri         & Arc          & Codegen    & Runtime config  & Ownership  & No$^*$ \\
Wails         & GC           & Reflection & None            & GC         & Yes \\
Neutralinojs  & Manual       & None       & Token auth      & Manual     & No \\
webview       & Manual       & None       & None            & Manual     & No \\
Gluon         & Ownership    & Serde      & None            & Ownership  & No \\
Sciter        & Refcount     & Typed      & None            & Mixed      & Partial \\
Ultralight    & Refcount     & None       & None            & Mixed      & Partial \\
Servo         & Ownership    & Typed      & Experimental    & Ownership  & No$^\dagger$ \\
Dioxus        & Ownership    & Typed      & None            & Ownership  & No \\
Flutter       & GC           & Codegen    & None            & GC         & Yes \\
\midrule
\textbf{Gossamer} & \textbf{Proved} & \textbf{Proved} & \textbf{Proved} & \textbf{Proved} & \textbf{Never} \\
\bottomrule
\end{tabular}
\par\smallskip
{\footnotesize $^*$Tauri uses \texttt{Arc<Mutex<...>>}, which is reference counting.
$^\dagger$Servo's content process runs a JavaScript GC.}
\end{table}

Table~\ref{tab:arch} compares architectural properties.

\begin{table}[ht]
\centering
\caption{Architecture comparison}
\label{tab:arch}
\small
\begin{tabular}{@{}lllcr@{}}
\toprule
\textbf{Framework} & \textbf{Backend} & \textbf{Rendering} & \textbf{Ships Browser?} & \textbf{Binary Size} \\
\midrule
Electron      & Node.js   & Chromium           & Yes  & 150--300\,MB \\
NW.js         & Node.js   & Chromium           & Yes  & 150--250\,MB \\
CEF           & C++       & Chromium           & Yes  & 100--200\,MB \\
Tauri         & Rust      & OS webview         & No   & 3--8\,MB \\
Wails         & Go        & OS webview         & No   & 5--10\,MB \\
Neutralinojs  & C++       & OS webview         & No   & 2--5\,MB \\
webview       & C         & OS webview         & No   & $<$1\,MB \\
Gluon         & Rust      & Installed browser  & No   & $<$1\,MB \\
Sciter        & C++       & Custom             & Own  & 5--8\,MB \\
Ultralight    & C++       & WebCore (stripped)  & Partial & 15--30\,MB \\
Servo         & Rust      & Servo              & Embedded & 30--50\,MB \\
Dioxus        & Rust      & Webview/native     & Optional & 3--10\,MB \\
Flutter       & Dart      & Skia               & No   & 10--20\,MB \\
\midrule
\textbf{Gossamer} & \textbf{Ephapax} & \textbf{OS webview} & \textbf{No} & \textbf{1--3\,MB} \\
\bottomrule
\end{tabular}
\end{table}

\section{Related Work}
\label{sec:related}

\paragraph{Linear types in systems programming.}
Rust's ownership system is affine (values may be used at most once),
with \texttt{Arc} and \texttt{Rc} providing shared ownership through
reference counting. Linear Haskell~\cite{bernardy2018} adds linearity
annotations to GHC but does not eliminate GC. The Clean programming
language uses uniqueness types (dual to linear types) for I/O.
Gossamer builds on Ephapax's dyadic approach, which offers both affine
and linear modes in a single language, and combines them with regions
to eliminate GC entirely.

\paragraph{Region-based memory management.}
MLKit~\cite{tofte1997} pioneered region inference for Standard ML.
Cyclone~\cite{jim2002cyclone} added regions to C with existential types.
Rust's lifetime system is region-adjacent but relies on borrow checking
rather than explicit regions. Ephapax's contribution is the combination
of explicit regions with linear types that provably prevent region
escape, eliminating the need for conservative inference or runtime
checks.

\paragraph{Typed IPC.}
Cap'n Proto~\cite{capnproto} and FlatBuffers~\cite{flatbuffers}
provide schema-defined serialisation with code generation, offering
partial type safety across process boundaries. gRPC provides typed
service definitions via Protocol Buffers. However, none of these
systems integrate with the type system of the host language to provide
compile-time guarantees --- they are external schemas with generated
code. Gossamer's IPC types are native Ephapax types checked by the
same compiler that checks the application logic.

\paragraph{Capability-based security.}
Tauri's capability system~\cite{tauri2024} is the most developed in the
webview shell space, but it operates at runtime. The E programming
language~\cite{miller2006} pioneered object-capability security.
Wasm-capabilities~\cite{wasi} provide a capability-based security model
for WebAssembly. Gossamer's contribution is encoding capabilities as
linear types, preventing duplication and enabling compile-time
enforcement.

\section{Future Work}
\label{sec:future}

\begin{itemize}[nosep]
  \item \textbf{Cross-platform:} extend Zig FFI to WKWebView (macOS)
    and WebView2 (Windows).
  \item \textbf{Mobile:} iOS and Android support via platform-specific
    FFI extensions.
  \item \textbf{Plugin system:} linear capability grants for third-party
    plugins with type-level isolation guarantees.
  \item \textbf{Frontend type bridge:} integrate with ReScript to
    provide end-to-end type safety from UI event to system call.
  \item \textbf{Formal verification:} extend Idris2 ABI proofs to cover
    the full IPC protocol, including serialisation correctness.
  \item \textbf{Benchmarks:} measure startup time, IPC latency, and
    memory usage against Tauri and Electron baselines.
\end{itemize}

\section{Conclusion}
\label{sec:conclusion}

Gossamer demonstrates that a webview shell framework can provide
compile-time proofs of resource safety, IPC type agreement, and
capability enforcement. By building on Ephapax's dyadic type system
(affine + linear modes) and region-based memory management, Gossamer
eliminates garbage collection entirely while guaranteeing that webview
handles cannot leak, IPC messages cannot mismatch, and permissions
cannot be bypassed.

The framework occupies a novel position in the design space: it is as
lightweight as the minimal C-based webview libraries ($\sim$1--3\,MB)
while providing stronger safety guarantees than any existing framework,
including Tauri's capability-based security model. The key insight is
that linear types and regions are not merely an alternative to garbage
collection --- they enable \emph{categorical} elimination of resource
management bugs that GC, reference counting, and manual tracking can
only \emph{mitigate}.

Gossamer is, to our knowledge, the first webview shell where the
compiler proves your application cannot leak resources, cannot use
freed handles, cannot bypass permissions, and will never require a
garbage collector.

\bibliographystyle{plain}

\begin{thebibliography}{99}

\bibitem{ephapax2025}
J.~D.~A. Jewell, ``Ephapax: a dyadic language with affine and linear
type modes,'' hyperpolymath, 2025.
\url{https://github.com/hyperpolymath/ephapax}

\bibitem{tofte1997}
M.~Tofte and J.-P.~Talpin, ``Region-based memory management,''
\textit{Information and Computation}, vol.~132, no.~2, pp.~109--176,
1997.

\bibitem{jim2002cyclone}
T.~Jim, J.~G. Morrisett, D.~Grossman, M.~W. Hicks, J.~Cheney, and
Y.~Wang, ``Cyclone: a safe dialect of C,'' in \textit{USENIX Annual
Technical Conference}, 2002.

\bibitem{bernardy2018}
J.-P. Bernardy, M.~Boespflug, R.~R. Newton, S.~Peyton Jones, and
A.~Spiwack, ``Linear Haskell: practical linearity in a higher-order
polymorphic language,'' \textit{Proceedings of the ACM on Programming
Languages}, vol.~2, no.~POPL, 2018.

\bibitem{wadler1990}
P.~Wadler, ``Linear types can change the world!'' in \textit{IFIP TC~2
Working Conference on Programming Concepts and Methods}, 1990.

\bibitem{walker2005}
D.~Walker, ``Substructural type systems,'' in \textit{Advanced Topics
in Types and Programming Languages} (B.~C. Pierce, ed.), MIT Press,
2005.

\bibitem{capnproto}
K.~Varda, ``Cap'n Proto,'' \url{https://capnproto.org/}, 2013.

\bibitem{flatbuffers}
W.~van Oortmerssen, ``FlatBuffers,''
\url{https://google.github.io/flatbuffers/}, 2014.

\bibitem{tauri2024}
Tauri Contributors, ``Tauri --- build smaller, faster, and more secure
desktop and mobile applications,'' \url{https://tauri.app/}, 2024.

\bibitem{miller2006}
M.~S. Miller, ``Robust composition: towards a unified approach to
access control and concurrency control,'' PhD thesis, Johns Hopkins
University, 2006.

\bibitem{wasi}
``WASI --- the WebAssembly System Interface,''
\url{https://wasi.dev/}, 2024.

\bibitem{hyperpolymath-abi}
J.~D.~A. Jewell, ``Idris2 ABI / Zig FFI universal standard,''
hyperpolymath, 2026.

\end{thebibliography}

\end{document}
